\chapter{Widgets And Layouts}

Widgets are the visible building blocks of a Qt application. Layouts decide how those widgets share space. This chapter introduces the core widget surface in \api{crystal-qt6}: common controls, sizing, attributes, layout composition, styling, and small bits of top-level window polish.

\section{The Widget Base}

\api{Qt6::Widget} wraps \api{QWidget}, the base class for most visible controls. It exposes shared behavior such as titles, visibility, enabled state, focus, tooltips, sizing, style sheets, cursor shape, drop acceptance, widget attributes, repaint scheduling, and screenshot capture.

\begin{lstlisting}[style=crystal]
panel = Qt6::Widget.new
panel.tool_tip = "Layer inspector"
panel.enabled = true
panel.set_minimum_size(280, 220)
panel.set_size_policy(Qt6::SizePolicy::Preferred, Qt6::SizePolicy::Expanding)
\end{lstlisting}

A widget can be top-level, parented to another widget, or placed inside a layout. Top-level widgets need a title and usually an initial size. Child widgets usually let their layout choose the final geometry.

\begin{lstlisting}[style=crystal]
window = Qt6::Widget.new
window.window_title = "Inspector"
window.resize(420, 320)
window.show
\end{lstlisting}

\section{Core Controls}

The simplest controls are labels and buttons. Labels show text. Buttons emit callbacks when clicked.

\begin{lstlisting}[style=crystal]
count = 0
value = Qt6::Label.new("0")
button = Qt6::PushButton.new("Increment")

button.on_clicked do
  count += 1
  value.text = count.to_s
end
\end{lstlisting}

Forms are usually built from a small family of input widgets:

\begin{lstlisting}[style=crystal]
name = Qt6::LineEdit.new("Terrain")
visible = Qt6::CheckBox.new("Visible")
kind = Qt6::ComboBox.new
kind << "Hexes" << "Terrain" << "Units"

opacity = Qt6::SpinBox.new
opacity.set_range(0, 100)
opacity.value = 72

scale = Qt6::DoubleSpinBox.new
scale.set_range(0.25, 4.0)
scale.value = 1.0
\end{lstlisting}

Qt forms often need a little desktop-style polish beyond the basic value editors. Editable combo boxes can accept freeform values while still offering known choices. Buttons can opt into dialog-default styling or flatter inspector-style presentation. Check boxes can expose the full Qt check-state enum when a partially checked state carries meaning.

\begin{lstlisting}[style=crystal]
kind.editable = true
kind.insert_item(1, "Overlay")
kind.current_text = "Terrain"

apply = Qt6::PushButton.new("Apply")
apply.default = true
apply.auto_default = false
apply.flat = false

visibility = Qt6::CheckBox.new("Visible")
visibility.tristate = true
visibility.check_state = Qt6::CheckState::PartiallyChecked

visibility.on_state_changed do |state|
  status.text = "Visibility state: #{state}"
end

mode = Qt6::RadioButton.new("Preview Mode")
mode.auto_exclusive = false
\end{lstlisting}

When several buttons represent a mode choice, group them with \api{ButtonGroup}. Groups can be exclusive like radio buttons, or they can simply provide shared callbacks and stable integer ids for toolbar-style toggle buttons.

\begin{lstlisting}[style=crystal]
terrain = Qt6::RadioButton.new("Terrain")
units = Qt6::RadioButton.new("Units")
roads = Qt6::RadioButton.new("Roads")

layer_modes = Qt6::ButtonGroup.new
layer_modes.add(terrain, 10)
layer_modes.add(units, 20)
layer_modes.add(roads, 30)
layer_modes.exclusive = true

layer_modes.on_id_clicked do |id|
  activate_layer_mode(id)
end

terrain.checked = true
\end{lstlisting}

When a text field should complete from a known set, attach a \api{Completer}. This is especially useful for search fields, tags, layer names, and command palettes.

\begin{lstlisting}[style=crystal]
search = Qt6::LineEdit.new
completer = Qt6::Completer.new(["Terrain", "Roads", "Rivers", "Routes"])
completer.wrap_around = true
completer.max_visible_items = 6
search.completer = completer
\end{lstlisting}

The \api{crystal-qt6} library also includes sliders, group boxes, progress bars, text editors, item widgets, date/time controls, tab widgets, scroll areas, splitters, and other common desktop controls. Later chapters use those in larger shell and editor examples.

Sliders are useful when the value should feel continuous. Besides \api{on_value_changed}, the binding also exposes handle press/release callbacks and an opt-in track-click mode that jumps directly to the clicked position.

\begin{lstlisting}[style=crystal]
zoom = Qt6::Slider.new(Qt6::Orientation::Horizontal)
zoom.set_range(25, 400)
zoom.value = 100
zoom.click_to_position = true

zoom.on_pressed do
  status.text = "Adjusting zoom"
end

zoom.on_value_changed do |value|
  preview_scale = value / 100.0
end

zoom.on_released do
  status.text = "Zoom #{zoom.value}%"
end
\end{lstlisting}

Date-oriented tools can move beyond spin boxes and text fields when the user needs to browse an actual month view. \api{CalendarWidget} works well for scheduling, journaling, replay controls, and any workflow where the visible page matters in addition to the selected date.

\begin{lstlisting}[style=crystal]
calendar = Qt6::CalendarWidget.new
calendar.set_date_range(
  Qt6::QDate.new(2026, 1, 1),
  Qt6::QDate.new(2026, 12, 31)
)
calendar.navigation_bar_visible = true
calendar.grid_visible = true
calendar.selected_date = Qt6::QDate.new(2026, 5, 5)

calendar.on_current_page_changed do |year, month|
  status.text = "Showing #{year}-#{month}"
end

calendar.on_clicked do |date|
  load_daily_summary(date)
end
\end{lstlisting}

\begin{figure}[htbp]
  \centering
  \screenshotimage{docs/book/images/widgets-core-controls-panel.png}{width=0.72\textwidth,height=0.62\textheight,keepaspectratio}%
  \caption{Core widget controls in one small preferences-style panel: line edits, combo boxes, spin boxes, sliders, radio buttons, and a calendar widget.}
\end{figure}

\section{Panel Polish Widgets}

Desktop editors often need a few small shell widgets in addition to the primary value editors. \api{ToolBox} is a compact stacked-page container that works well for inspectors and sidebars. \api{FocusFrame} lets Qt draw a native-looking focus highlight around another widget. \api{SizeGrip} provides an explicit resize corner for custom dialogs or utility windows that do not already expose one through a status bar.

\begin{lstlisting}[style=crystal]
toolbox = Qt6::ToolBox.new
toolbox.add_item(Qt6::Label.new("Layer list"), "Layers")
toolbox.add_item(Qt6::Label.new("Brush settings"), "Brushes")
toolbox.current_index = 1

search = Qt6::LineEdit.new("terrain")
focus_frame = Qt6::FocusFrame.new
focus_frame.widget = search

corner_grip = Qt6::SizeGrip.new(window)
\end{lstlisting}

Use \api{ToolBox} when a tab bar would feel too heavy and only one inspector page should be visible at a time. Use \api{FocusFrame} sparingly, usually when the focused editor should stand out inside a busy shell. Use \api{SizeGrip} when the platform or window style does not already provide an obvious resize affordance.

\section{Common Control Choices}

Most forms become easier to design when you choose the narrowest control that fits the value being edited.

\begin{longtable}{@{}p{0.30\textwidth}p{0.64\textwidth}@{}}
\toprule
Control & Good fit \\
\midrule
\api{Label} & Static or read-only text. \\
\api{PushButton} & Direct commands such as Add, Reset, Export, or Apply. \\
\api{LineEdit} & Short single-line text such as a name, path, or search term. \\
\api{PlainTextEdit} & Multiline notes, logs, source snippets, and longer plain text. \\
\api{CheckBox} & Independent boolean or tri-state options. \\
\api{RadioButton} & One choice from a small visible set, often paired with \api{ButtonGroup}. \\
\api{ComboBox} & A compact choice from a known option list. \\
\api{SpinBox} and \api{DoubleSpinBox} & Bounded numeric values. \\
\api{Slider} and \api{Dial} & Approximate numeric values where immediate feedback matters. \\
\api{CalendarWidget} & Visual month-based date picking and schedule browsing. \\
\api{ProgressBar} & Progress, completion, or a read-only bounded value. \\
Item widgets & \api{ListWidget}, \api{TreeWidget}, and \api{TableWidget} are good for small owned item collections. For larger or shared data, prefer the model/view APIs in Chapter 10. \\
\bottomrule
\end{longtable}

\section{Visibility, Enabled State, And Focus}

Visibility and enabled state are shared widget features:

\begin{lstlisting}[style=crystal]
advanced_panel.visible = show_advanced.checked?
apply_button.enabled = !name.text.empty?
\end{lstlisting}

Focus is explicit when you need keyboard input to land in a particular control.

\begin{lstlisting}[style=crystal]
name.focus_policy = Qt6::FocusPolicy::StrongFocus
name.set_focus
\end{lstlisting}

For ordinary forms, Qt's default focus behavior is usually enough. Use explicit focus control for editor workflows, search fields, modal flows, and custom canvases.

\section{Sizing}

Qt layouts negotiate size from several pieces of information: each widget's size hint, its minimum and maximum size, and its size policy. \api{crystal-qt6} exposes the common controls for that negotiation.

\begin{lstlisting}[style=crystal]
preview = Qt6::Label.new("Preview")
preview.set_minimum_size(240, 160)
preview.maximum_height = 240

swatch = Qt6::Label.new("")
swatch.set_fixed_size(96, 24)

notes = Qt6::PlainTextEdit.new
notes.set_size_policy(Qt6::SizePolicy::Expanding, Qt6::SizePolicy::Expanding)
\end{lstlisting}

Use fixed sizes sparingly. They are appropriate for swatches, icon previews, compact counters, and controls whose physical dimensions are part of the UI. For panels and editors, prefer minimum sizes plus expanding or preferred size policies.

\section{Widget Attributes And Cursors}

Qt widget attributes are lower-level flags for behavior that is not always worth a dedicated high-level method. \api{crystal-qt6} exposes them through \api{WidgetAttribute}.

\begin{lstlisting}[style=crystal]
overlay = Qt6::Label.new("Drop files here")
overlay.set_attribute(Qt6::WidgetAttribute::StyledBackground)
overlay.set_attribute(Qt6::WidgetAttribute::TransparentForMouseEvents)
\end{lstlisting}

Some attributes also have convenience methods. Prefer the clearer method when one exists:

\begin{lstlisting}[style=crystal]
canvas = Qt6::EventWidget.new
canvas.mouse_tracking = true
canvas.accept_drops = true
canvas.cursor_shape = Qt6::CursorShape::Cross
\end{lstlisting}

Attributes are powerful, but they are close to Qt itself. Use them to express known Qt behavior, not as a substitute for ordinary widget APIs.

\section{Layout Families}

Qt layouts are parent-owned objects that place child widgets. In Crystal, \api{Widget#vbox}, \api{#hbox}, \api{#form}, and \api{#grid} create the corresponding layout and yield it for composition.

\api{VBoxLayout} stacks widgets vertically:

\begin{lstlisting}[style=crystal]
panel = Qt6::Widget.new
panel.vbox do |column|
  column << Qt6::Label.new("Layers")
  column << layer_list
  column << apply_button
end
\end{lstlisting}

\api{HBoxLayout} is useful for command rows:

\begin{lstlisting}[style=crystal]
buttons = Qt6::Widget.new
buttons.hbox do |row|
  row << Qt6::PushButton.new("Add")
  row << Qt6::PushButton.new("Remove")
  row << Qt6::PushButton.new("Apply")
end
\end{lstlisting}

Use \api{add_stretch} when one part of the row or column should absorb the extra space instead of forcing widgets apart evenly:

\begin{lstlisting}[style=crystal]
buttons = Qt6::Widget.new
buttons.hbox do |row|
  row << Qt6::PushButton.new("Add")
  row << Qt6::PushButton.new("Remove")
  row.add_stretch
  row << Qt6::PushButton.new("Apply")
end
\end{lstlisting}

The same idea works in a vertical stack when a status row or tool strip should stay pinned to one edge:

\begin{lstlisting}[style=crystal]
panel = Qt6::Widget.new
panel.vbox do |column|
  column << layer_list
  column.add_stretch
  column << status_bar
end
\end{lstlisting}

\api{FormLayout} is best for labeled editor fields:

\begin{lstlisting}[style=crystal]
properties = Qt6::Widget.new
properties.form do |form|
  form.add_row("Name", name)
  form.add_row("Kind", kind)
  form.add_row("Opacity", opacity)
  form.add_row(visible)
end
\end{lstlisting}

\api{GridLayout} works well for matrix-like tools, command palettes, and compact dashboards:

\begin{lstlisting}[style=crystal]
grid = Qt6::Widget.new
grid.grid do |layout|
  layout.add(Qt6::Label.new("X"), 0, 0)
  layout.add(Qt6::Label.new("Y"), 0, 1)
  layout.add(x_spin, 1, 0)
  layout.add(y_spin, 1, 1)
end
\end{lstlisting}

\section{Container Widgets}

Container widgets give a layout more structure than rows and columns alone. A \api{GroupBox} labels related controls and can optionally become checkable when the whole group represents an enabled feature.

\begin{lstlisting}[style=crystal]
appearance = Qt6::GroupBox.new("Appearance")
appearance.checkable = true
appearance.checked = true
appearance.alignment = Qt6::AlignmentFlag::Center
appearance.flat = true
appearance.form do |form|
  form.add_row("Opacity", opacity)
  form.add_row("Scale", scale)
end
\end{lstlisting}

\api{TabWidget} is useful when several peer pages need the same space. Keep tab labels short, and avoid hiding primary workflow controls several tabs deep.

\begin{lstlisting}[style=crystal]
tabs = Qt6::TabWidget.new
tabs.add_tab(general_page, "General")
tabs.add_tab(editor_page, "Editor")
tabs.set_tab_text(1, "Editing")
tabs.current_widget = general_page
\end{lstlisting}

Use \api{ScrollArea} when a panel can legitimately grow beyond the available window height. The common form pattern is to place one content widget inside the scroll area and let the area resize it with the viewport.

\begin{lstlisting}[style=crystal]
scroll = Qt6::ScrollArea.new
scroll.widget = settings_panel
scroll.widget_resizable = true
scroll.alignment = Qt6::AlignmentFlag::Top
\end{lstlisting}

\api{Splitter} gives the user direct control over neighboring regions, such as a layer list and a property inspector.

\begin{lstlisting}[style=crystal]
splitter = Qt6::Splitter.new(Qt6::Orientation::Horizontal)
splitter << layer_list
splitter << properties_panel
splitter.children_collapsible = false
splitter.set_sizes([280, 520])
\end{lstlisting}

\api{StackedWidget} is for mode switching when only one page should be visible and some other control chooses the active page.

\begin{lstlisting}[style=crystal]
stack = Qt6::StackedWidget.new
stack << simple_page
stack << advanced_page
mode_combo.on_current_index_changed { |index| stack.current_index = index }
stack.current_widget = simple_page
\end{lstlisting}

\begin{figure}[htbp]
  \centering
  \screenshotimage{docs/book/images/widgets-container-workbench.png}{width=0.92\textwidth,height=0.44\textheight,keepaspectratio}%
  \caption{Container widgets working together: a \api{ToolBox} beside a \api{TabWidget}, with a \api{ScrollArea}, \api{Splitter}, and \api{StackedWidget} inside the workbench.}
\end{figure}

\section{Composing A Panel}

Most desktop panels combine several layouts. A common pattern is a vertical column containing a form plus a row of actions.

\begin{lstlisting}[style=crystal]
name = Qt6::LineEdit.new("Terrain")
kind = Qt6::ComboBox.new
kind << "Hexes" << "Terrain" << "Units"

visible = Qt6::CheckBox.new("Visible")
apply = Qt6::PushButton.new("Apply")
reset = Qt6::PushButton.new("Reset")

panel = Qt6::Widget.new
panel.set_minimum_size(320, 220)
panel.vbox do |column|
  form_host = Qt6::Widget.new
  form_host.form do |form|
    form.add_row("Name", name)
    form.add_row("Kind", kind)
    form.add_row(visible)
  end

  actions = Qt6::Widget.new
  actions.hbox do |row|
    row << reset
    row << apply
  end

  column << form_host
  column << actions
end
\end{lstlisting}

This is not the only way to compose a panel, but it scales well. The form stays responsible for labels and fields. The horizontal layout stays responsible for commands. The outer vertical layout controls overall flow.

\begin{figure}[H]
  \centering
  {\setlength{\fboxsep}{10pt}%
   \colorbox{black!10}{%
     \screenshotimage{docs/book/images/widgets-layer-inspector-panel.png}{width=0.58\textwidth,height=0.42\textheight,keepaspectratio}%
   }}%
  \caption{A small layer inspector built from form rows, grouped controls, and a command row.}
\end{figure}

\section{Layout Pitfalls}

Fixed geometry is the most common way to make a Qt interface brittle. Use fixed sizes for things whose dimensions are part of their meaning, such as swatches, icon previews, and compact counters. For editors and panels, prefer minimum sizes, maximum sizes, stretch, and size policies.

Avoid manually positioning child widgets inside a widget that already has a layout. Once a layout owns the geometry, let it do the placement. If a page needs a custom drawing surface, keep that as a dedicated canvas widget and place it beside ordinary controls with layouts or splitters.

Let parent and layout ownership do its job. In normal UI code, create widgets, add them to layouts or containers, and release only top-level objects or explicitly detached wrappers. If a panel becomes too tall, reach for \api{ScrollArea}, tabs, or a splitter before hard-coding coordinates or shrinking controls until labels no longer fit.

\section{Styling}

Qt style sheets are CSS-like strings applied to widgets or the whole application. They are useful for lightweight polish, warnings, previews, and demo styling.

\begin{lstlisting}[style=crystal]
swatch = Qt6::Label.new("")
swatch.set_fixed_size(96, 24)
swatch.style_sheet = "background-color: #3470b0; border: 1px solid #6b7280;"

name.tool_tip = "Layer name shown in the layer list"
\end{lstlisting}

Application-wide style sheets are available from the application wrapper:

\begin{lstlisting}[style=crystal]
app = Qt6.application
app.style_sheet = "
  QGroupBox { font-weight: bold; }
  QPushButton { padding: 4px 10px; }
"
\end{lstlisting}

Use style sheets to support the interface rather than replace layout. Sizing, hierarchy, and interaction should come from the widget tree and layout structure first.

\section{Window Polish}

Top-level widgets and windows expose title, icon, size, visibility, and capture helpers. Small display widgets can also add focused presentation polish. For example, \api{LcdNumber} now exposes decimal-point and overflow helpers that are useful for compact counters or telemetry panes.

\begin{lstlisting}[style=crystal]
main = Qt6::MainWindow.new
main.window_title = "Map Tool"
main.resize(960, 640)

icon = Qt6::QIcon.new("assets/app-icon.png")
main.window_icon = icon

main.status_bar.show_message("Ready")
main.show
\end{lstlisting}

The application object can also hold application-wide metadata and window icon state:

\begin{lstlisting}[style=crystal]
app = Qt6.application
app.name = "Map Tool"
app.organization_name = "Example Studio"
app.window_icon = icon
\end{lstlisting}

\begin{lstlisting}[style=crystal]
fps = Qt6::LcdNumber.new
fps.digit_count = 5
fps.small_decimal_point = true
fps.display(59.9)
\end{lstlisting}

For quick exports and documentation screenshots, \api{Widget#grab} captures a widget into a \api{QPixmap}.

\begin{lstlisting}[style=crystal]
pixmap = main.grab
pixmap.save("window.png")
\end{lstlisting}

\section{Exercise: Build A Preferences Panel}

As a small practice project, build a preferences panel with two tabs. The General tab can hold a \api{LineEdit} for a profile name, a \api{ComboBox} for a theme choice, and a \api{CheckBox} for an option such as restoring the last project. The Editor tab can hold a \api{SpinBox} for tab width, a \api{Slider} for zoom, and a \api{PlainTextEdit} for notes.

Use a \api{FormLayout} for each page, put the pages in a \api{TabWidget}, and add Reset and Apply buttons below the tabs. If the panel grows, put the tab widget or its pages in a \api{ScrollArea}. For extra practice, call \api{Widget#grab} and save a screenshot of the result.

\section{Where To Look Next}

\api{examples/inspector_workbench.cr} is the best example for form controls, group boxes, sliders, spin boxes, tabs, splitters, and a live custom canvas. \api{examples/editor_shell.cr} shows how ordinary widgets fit into a main-window shell with actions, menus, toolbars, docks, dialogs, and a status bar.
